\documentclass[11pt,a4paper]{article}

\usepackage[a4paper,margin=1in]{geometry}
\usepackage{fontspec}
\IfFontExistsTF{Times New Roman}
  {\setmainfont{Times New Roman}}
  {\IfFontExistsTF{TeX Gyre Termes}{\setmainfont{TeX Gyre Termes}}{\setmainfont{Times}}}
\usepackage{newtxmath}
\usepackage{amsmath}
\usepackage{booktabs}
\usepackage{longtable}
\usepackage{array}
\usepackage{enumitem}
\usepackage{xcolor}
\usepackage{listings}
\usepackage{graphicx}
\usepackage{float}
\usepackage[round,authoryear]{natbib}
\usepackage[hidelinks]{hyperref}
\usepackage{xurl}
\usepackage{fancyhdr}
\usepackage{microtype}

\setlength{\parindent}{0pt}
\setlength{\parskip}{0.65em}
\setlength{\emergencystretch}{3em}
\setlist[itemize]{leftmargin=1.4em,itemsep=0.2em,topsep=0.2em}
\setlist[enumerate]{leftmargin=1.6em,itemsep=0.2em,topsep=0.2em}

\pagestyle{fancy}
\fancyhf{}
\lhead{\small Ticket 06 Study Report}
\rhead{\small No-Dividend American Call Validation}
\cfoot{\thepage}
\renewcommand{\headrulewidth}{0.3pt}

\lstset{
  basicstyle=\ttfamily\small,
  breaklines=true,
  frame=single,
  columns=fullflexible,
  keepspaces=true,
  rulecolor=\color{black!30},
  xleftmargin=0.5em,
  xrightmargin=0.5em
}

\newcommand{\code}[1]{\texttt{#1}}
\newcommand{\file}[1]{\path{#1}}
\newcommand{\func}[1]{\path{#1}}
\newcolumntype{L}[1]{>{\raggedright\arraybackslash}p{#1}}
\newcolumntype{P}[1]{>{\raggedright\arraybackslash\ttfamily\footnotesize}p{#1}}

\title{\textbf{Ticket 06 Study Report: No-Dividend American Call Validation}\\
\large SURF2026 American Option Risk Surfaces}
\author{Codex-assisted implementation report}
\date{June 16, 2026}

\begin{document}
\maketitle

\begin{abstract}
Ticket 06 validates a special but important American option theorem: under the standard
Black-Scholes assumptions with no dividends, a standard American call should match the
corresponding European call. This ticket uses the Ticket 04 American CN/PSOR solver as-is,
sets \(q=0\), compares against the Ticket 01 European Black-Scholes call utility, records
obstacle and PSOR convergence diagnostics, and creates two validation figures. The purpose is
to check that the American exercise machinery does not create artificial early-exercise value.
This report is written for beginner researchers with accounting, data, machine-learning, and
Python experience who are still learning American option pricing and numerical validation.
\end{abstract}

\tableofcontents
\newpage

\section{Role of Ticket 06 in the Whole Solver Project}

The project is building a validated classical CN/PSOR benchmark before moving toward boundary
diagnostics, Greeks, stress regimes, datasets, or neural surrogate models. The previous solver
tickets created the layers needed for this validation ladder:
\begin{itemize}
  \item Ticket 01 implemented payoff functions and European Black-Scholes closed-form utilities.
  \item Ticket 02 implemented uniform grids, finite-difference operator helpers, and boundary
        value helpers.
  \item Ticket 03 validated European Crank-Nicolson time stepping against closed-form prices.
  \item Ticket 04 implemented the reusable American CN/PSOR obstacle core.
  \item Ticket 05 validated American put behavior in a representative early-exercise case.
\end{itemize}

Ticket 06 checks a different American option behavior. It asks whether the same American
CN/PSOR machinery avoids inventing early exercise when theory says early exercise should not be
valuable. This is a validation ticket, not a new solver ticket. No core solver source file is
changed, and no free-boundary extraction, Greeks, stress maps, datasets, or neural-network files
are introduced.

\section{Theory and Methodology/Literature Connection}

The documented source basis for Ticket 06 is internal:
\begin{itemize}
  \item \file{reports/01_literature/literature_map.md};
  \item \file{reports/02_math/formulation_note.md};
  \item \file{reports/03_solver/solver_validation_plan.md};
  \item \file{reports/03_solver/solver_implementation_tickets.md};
  \item the Ticket 01 through Ticket 05 reports and code.
\end{itemize}

The literature map identifies the no-dividend American call theorem as a key validation check.
The formulation note states that when \(q=0\), a standard American call should match the
European call under the usual Black-Scholes assumptions. The solver validation plan places this
check after European closed-form validation and before the later dividend-paying call work.

This reference integration pass adds verified citations for the no-dividend American call control
logic. The option-pricing framework follows Black-Scholes and Merton
\citep{black_scholes_1973,merton_1973}, while the American exercise machinery remains grounded in
the project's LCP/PSOR validation ladder and standard LCP methodology
\citep{cottle_pang_stone_2009}.

\section{No-Dividend American Call Theorem in Plain Language}

A call option gives the holder the right to buy the underlying asset for strike \(K\). Its payoff
is
\begin{equation}
  \Phi_{\mathrm{call}}(S) = \max(S-K,0).
\end{equation}

An American call can be exercised before maturity. A European call can be exercised only at
maturity. At first glance, the American call seems like it must be more valuable because it gives
the holder more flexibility.

For a standard call with no dividends, however, early exercise is not optimal under the usual
Black-Scholes assumptions. Exercising early gives up remaining optionality. It also requires
paying the strike earlier than necessary. With \(q=0\), there is no dividend benefit from owning
the stock earlier. Therefore, the extra American exercise right should have no positive value:
\begin{equation}
  C_{\mathrm{American}}(S; q=0) = C_{\mathrm{European}}(S; q=0).
  \label{eq:no-dividend-call-theorem}
\end{equation}

This is not saying American calls are always equal to European calls. Dividend-paying calls are
different because dividends can make early stock ownership valuable. Ticket 06 is only the
no-dividend control case.

\section{Why Ticket 06 Follows American Put Validation}

Ticket 05 tested the American put case, where early exercise can be financially meaningful. The
American put should dominate payoff and can exceed the European put. Ticket 06 tests the opposite
kind of sanity check: an American option path is active, PSOR projection is active, and the payoff
obstacle is present, but theory says the exercise feature should not create genuine value.

This is a powerful check because a solver can pass \(U \geq \Phi\) and still be wrong. If a
no-dividend American call is meaningfully above the European call, the solver may be creating
artificial early exercise. That could come from a wrong call payoff, a wrong call boundary value,
a hidden nonzero dividend yield, a bad boundary adjustment, or projection logic that distorts the
continuation solution.

\section{Validation Criteria}

\subsection{American Equals European When q = 0}

The main Ticket 06 criterion is approximate equality:
\begin{equation}
  U_{\mathrm{American\ call}}(S,T; q=0)
  \approx C_{\mathrm{European}}(S,T; q=0).
\end{equation}

The experiment reports the maximum absolute error and RMSE over the target moneyness region
\[
  0.4 \leq S/K \leq 1.8.
\]
This target region follows the solver validation plan. It focuses attention on the region that is
likely to matter later for risk surfaces.

\subsection{Obstacle Condition}

Even though early exercise should not be optimal, the American solver must still respect the
payoff obstacle:
\begin{equation}
  U(S,\tau) \geq \Phi_{\mathrm{call}}(S)
\end{equation}
at all grid nodes and time levels. Ticket 06 reports the minimum obstacle gap and maximum
obstacle violation.

\subsection{No False Early-Exercise Premium}

The American-minus-European difference is interpreted as a no-false-early-exercise diagnostic:
\begin{equation}
  U_{\mathrm{American\ call}}(S,T) - C_{\mathrm{European}}(S,T).
\end{equation}
For \(q=0\), this should remain close to zero. Small positive or negative values can occur from
finite-difference and boundary approximation error. The report does not treat roundoff-scale or
small discretization-scale differences as financial early exercise.

\subsection{PSOR Convergence Metadata}

Every time step returns PSOR metadata. Ticket 06 reports whether all steps converged, how many
steps were solved, the maximum and mean iteration counts, and the maximum final update. This
matters because a price comparison is not credible if the LCP solves did not converge.

\section{What Ticket 06 Is Not}

Ticket 06 is not dividend-paying American call validation. It uses \(q=0\) only. Positive
dividend yields belong to a later ticket because they can make early exercise meaningful at high
spot values.

Ticket 06 is also not boundary extraction. It does not compute a continuation premium curve, does
not locate an exercise boundary, and does not label any figure with a boundary. The theorem-error
figure shows American-minus-European pricing error only.

\section{Implementation Design Decisions}

\begin{longtable}{L{0.28\textwidth}L{0.60\textwidth}}
\toprule
Decision & Reason \\
\midrule
\endhead
Use the Ticket 04 solver as-is &
Ticket 06 is a validation layer. Solver edits are out of scope unless a genuine bug is found. \\
Set \(q=0.0\) in every case &
The no-dividend theorem only applies to standard no-dividend calls. \\
Use \(K=1\), \(S_{\max}=4K\), and \(S/K\in[0.4,1.8]\) &
These follow the normalized-strike workflow and target region used in earlier tickets. \\
Use medium and fine grids &
This gives a light refinement check without turning Ticket 06 into a full grid/domain sensitivity
study. \\
Write CSV tables &
Tables make the validation evidence reproducible and easy to inspect. \\
Generate two figures &
The value comparison shows overlap visually; the theorem-error figure shows the small numerical
difference directly. \\
\bottomrule
\end{longtable}

\section{Code-Level Function Explanations}

{\small
\begin{longtable}{L{0.38\textwidth}L{0.50\textwidth}}
\toprule
Function & Role \\
\midrule
\endhead
\func{run_validation_cases} &
Runs the medium and fine no-dividend American call cases, builds validation rows, selected-spot
rows, and figure metadata. \\
\func{summarize_no_dividend_call_case} &
Computes theorem-error metrics, obstacle diagnostics, and PSOR convergence summaries for one
solver result. \\
\func{selected_spot_rows} &
Reports payoff, European call, American call, American-minus-European, and American-minus-payoff
at selected moneyness points. \\
\func{nearest_spot_index} &
Chooses the nearest grid node to a target spot so selected-spot rows are reproducible. \\
\func{write_csv} &
Writes stable CSV artifacts with fixed column ordering. \\
\func{create_value_comparison_figure} &
Creates the American call, European call, and payoff comparison figure if matplotlib is available. \\
\func{create_theorem_error_figure} &
Creates the American-minus-European error figure if matplotlib is available. \\
\code{main} &
Runs the full Ticket 06 experiment and writes CSV and figure artifacts. \\
\bottomrule
\end{longtable}
}

\section{Validation Result Tables and Interpretation}

\subsection{Validation Cases}

{\small
\begin{longtable}{@{}L{0.16\textwidth}L{0.07\textwidth}L{0.07\textwidth}L{0.07\textwidth}L{0.07\textwidth}L{0.09\textwidth}L{0.07\textwidth}L{0.07\textwidth}@{}}
\toprule
Case & \(K\) & \(T\) & \(r\) & \(q\) & \(\sigma\) & \(M\) & \(N\) \\
\midrule
\endhead
Medium & 1.0 & 1.0 & 0.05 & 0.00 & 0.20 & 120 & 120 \\
Fine & 1.0 & 1.0 & 0.05 & 0.00 & 0.20 & 180 & 180 \\
\bottomrule
\end{longtable}
}

\subsection{American-Versus-European Error Summary}

{\small
\begin{longtable}{@{}L{0.13\textwidth}L{0.16\textwidth}L{0.14\textwidth}L{0.14\textwidth}L{0.14\textwidth}L{0.12\textwidth}@{}}
\toprule
Case & Max abs error & RMSE & Min Am. - Eur. & Max Am. - Eur. & Max-error spot \\
\midrule
\endhead
Medium & \(2.7817\times10^{-4}\) & \(1.1977\times10^{-4}\) & \(-2.7817\times10^{-4}\) & \(4.3704\times10^{-5}\) & 0.966667 \\
Fine & \(1.2357\times10^{-4}\) & \(5.3375\times10^{-5}\) & \(-1.2357\times10^{-4}\) & \(1.9771\times10^{-5}\) & 0.977778 \\
\bottomrule
\end{longtable}
}

The fine-grid maximum absolute error and RMSE are smaller than the medium-grid values. The
American-minus-European differences are small finite-difference errors rather than evidence of
meaningful early exercise. The maximum positive American-minus-European value is below
\(4.4\times10^{-5}\) on the medium grid and below \(2.0\times10^{-5}\) on the fine grid.

\subsection{Obstacle Summary}

\begin{longtable}{@{}L{0.16\textwidth}L{0.18\textwidth}L{0.22\textwidth}@{}}
\toprule
Case & Min obstacle gap & Max obstacle violation \\
\midrule
\endhead
Medium & 0 & 0 \\
Fine & 0 & 0 \\
\bottomrule
\end{longtable}

The obstacle diagnostics show no violation in the generated value grids. This confirms that the
American solver respected \(U\geq\Phi\), but it does not by itself prove the no-dividend theorem.
That is why the European comparison is the central Ticket 06 check.

\subsection{PSOR Convergence Metadata}

\begin{longtable}{@{}L{0.14\textwidth}L{0.18\textwidth}L{0.15\textwidth}L{0.18\textwidth}L{0.22\textwidth}@{}}
\toprule
Case & All steps converged & Step count & Max iterations & Max final update \\
\midrule
\endhead
Medium & True & 120 & 13 & \(9.2800\times10^{-9}\) \\
Fine & True & 180 & 15 & \(7.3158\times10^{-9}\) \\
\bottomrule
\end{longtable}

Every PSOR step converged in both validation cases. The final updates are below the intended
tolerance scale, so the reported price comparisons are not hiding a visible PSOR convergence
failure.

\subsection{Selected Spot Price Comparison}

{\footnotesize
\begin{longtable}{@{}L{0.10\textwidth}L{0.10\textwidth}L{0.12\textwidth}L{0.13\textwidth}L{0.13\textwidth}L{0.13\textwidth}L{0.13\textwidth}@{}}
\toprule
Case & \(S/K\) & Payoff & European call & American call & Am. - Eur. & Am. - payoff \\
\midrule
\endhead
Medium & 0.5 & 0 & 0.000024 & 0.000032 & 0.000008 & 0.000032 \\
Medium & 0.8 & 0 & 0.018594 & 0.018498 & \(-0.000096\) & 0.018498 \\
Medium & 1.0 & 0 & 0.104506 & 0.104230 & \(-0.000276\) & 0.104230 \\
Medium & 1.2 & 0.200000 & 0.261690 & 0.261573 & \(-0.000118\) & 0.061573 \\
Medium & 1.5 & 0.500000 & 0.549701 & 0.549700 & \(-0.000001\) & 0.049700 \\
Fine & 0.488889 & 0 & 0.000016 & 0.000018 & 0.000003 & 0.000018 \\
Fine & 0.8 & 0 & 0.018594 & 0.018551 & \(-0.000043\) & 0.018551 \\
Fine & 1.0 & 0 & 0.104506 & 0.104384 & \(-0.000122\) & 0.104384 \\
Fine & 1.2 & 0.200000 & 0.261690 & 0.261638 & \(-0.000052\) & 0.061638 \\
Fine & 1.488889 & 0.488889 & 0.538692 & 0.538691 & \(-0.000001\) & 0.049802 \\
\bottomrule
\end{longtable}
}

The selected spot rows make the theorem check easier to read. Around the money, the American and
European call values are very close. The American-minus-payoff column remains positive because a
call can have time value even when early exercise is not optimal.

\subsection{Medium-to-Fine Comparison}

The largest selected-spot medium-to-fine value difference reported by the experiment is
\(0.0110089\). This number is driven by nearest-grid selection near \(S/K=1.5\): the medium grid
uses exactly \(S=1.5\), while the fine grid's nearest node is \(S=1.4888889\). It is useful as a
simple audit value, but it is not the same thing as an interpolated convergence study.

\section{Figure Interpretation}

Figure~\ref{fig:value-comparison} plots the fine-grid American call value, European call value,
and payoff over moneyness. The American and European value curves visually overlap in the target
region. The payoff line sits below the option value in the continuation-value region, which is
normal for a call with time remaining.

\begin{figure}[H]
\centering
\includegraphics[width=0.82\textwidth]{results/01_solver_validation/figures/ticket_06_no_dividend_american_call_value_comparison.png}
\caption{Ticket 06 no-dividend American call value, European call value, and payoff comparison.}
\label{fig:value-comparison}
\end{figure}

Figure~\ref{fig:theorem-error} plots American call value minus European call value. The error is
small, crosses near zero, and remains at the finite-difference scale reported in the CSV table.
This figure visualizes theorem error only. It does not mark or infer any exercise boundary.

\begin{figure}[H]
\centering
\includegraphics[width=0.82\textwidth]{results/01_solver_validation/figures/ticket_06_no_dividend_american_call_theorem_error.png}
\caption{Ticket 06 American-minus-European call error for the no-dividend theorem check.}
\label{fig:theorem-error}
\end{figure}

\section{Testing Strategy}

\begin{longtable}{L{0.27\textwidth}L{0.31\textwidth}L{0.30\textwidth}}
\toprule
Test category & What it proves & What it does not prove \\
\midrule
\endhead
Script/artifact tests &
The experiment runs and writes expected CSV and PNG artifacts. &
Full validation across all market regimes. \\
CSV schema tests &
The result tables have stable columns for later reporting. &
That every future diagnostic has already been designed. \\
\(q=0\) tests &
Every validation row uses the no-dividend setting. &
Dividend-paying call behavior. \\
Obstacle tests &
The American value does not fall below payoff within tolerance. &
That early exercise is economically present. \\
European equality tests &
American call values are close to European Black-Scholes calls. &
A formal convergence proof. \\
No-false-premium tests &
Positive American-minus-European values remain small. &
Positive-dividend call validation. \\
PSOR metadata tests &
All default time steps converged and expose iteration/update summaries. &
Convergence for every possible grid or relaxation parameter. \\
Nearest-grid test &
Selected spot reporting is predictable. &
Interpolation accuracy. \\
Figure tests &
Meaningful PNGs are created when matplotlib is available. &
That figures alone validate the solver. \\
Scope and docstring tests &
Ticket 06 files identify their ticket and do not introduce forbidden APIs/files. &
Future tickets that intentionally add diagnostics. \\
\bottomrule
\end{longtable}

\section{Test Results and Interpretation}

The test command was:
\begin{lstlisting}[language=bash]
PYTHONPYCACHEPREFIX=/private/tmp/codex_pycache PYTHONPATH=src MPLCONFIGDIR=/private/tmp/matplotlib-cache FC_CACHEDIR=/private/tmp/fontconfig-cache python3 -m unittest discover -s tests
\end{lstlisting}

Observed result during Ticket 06 implementation:
\begin{lstlisting}
Ran 63 tests

OK
\end{lstlisting}

The experiment command was:
\begin{lstlisting}[language=bash]
PYTHONPYCACHEPREFIX=/private/tmp/codex_pycache PYTHONPATH=src MPLCONFIGDIR=/private/tmp/matplotlib-cache FC_CACHEDIR=/private/tmp/fontconfig-cache python3 experiments/03_no_dividend_american_call_validation.py
\end{lstlisting}

The experiment wrote two validation rows, ten selected-spot rows, and two PNG figures.

\section{Implementation Pitfalls and Debugging Notes}

{\small
\begin{longtable}{@{}L{0.22\textwidth}L{0.24\textwidth}L{0.22\textwidth}L{0.18\textwidth}@{}}
\toprule
Recorded item or pitfall & Symptom & Fix or guard & Beginner lesson \\
\midrule
\endhead
Expected TDD red run &
The new tests failed with \code{FileNotFoundError} because the Ticket 06 experiment script did
not exist yet. &
Create \file{experiments/03_no_dividend_american_call_validation.py}. &
A red test confirms the test is aimed at the new artifact. \\
Nearest-grid threshold too strict &
The selected medium-to-fine value difference was \(0.0110089\), just above an initial \(0.01\)
test threshold. &
Keep the theorem-error improvement strict, but allow a modest nearest-node selected-spot
difference below \(0.012\). &
Nearest-node tables are audits, not interpolated convergence studies. \\
Accidentally using \(q>0\) &
The theorem would no longer apply. &
Validate every Ticket 06 case and row has \(q=0.0\). &
Always check assumptions that define a theorem. \\
Overclaiming exact equality &
Finite-difference values can be slightly above or below closed form. &
Report max absolute error and RMSE instead of claiming exact equality. &
Numerical validation is tolerance-based. \\
Confusing this with dividend-call validation &
Positive dividends can change the financial conclusion. &
State clearly that positive \(q\) belongs to a later ticket. &
Control cases and target cases answer different questions. \\
Ignoring PSOR metadata &
A small price error could hide failed iterations. &
Report convergence status, iteration counts, and final update values. &
Numerical credibility includes solver-status evidence. \\
\bottomrule
\end{longtable}
}

\section{Limitations}

\begin{longtable}{L{0.34\textwidth}L{0.54\textwidth}}
\toprule
Not included in Ticket 06 & Reason \\
\midrule
\endhead
Dividend-paying American call validation & Positive-dividend calls belong to a later ticket. \\
American put validation & Already covered in Ticket 05. \\
Free-boundary extraction & Ticket 06 does not infer or label exercise boundaries. \\
Continuation-premium module & Boundary-oriented diagnostics are later work. \\
Greeks & Delta and Gamma diagnostics are later work. \\
Stress maps, datasets, neural networks & These require a more complete validated solver pipeline. \\
Full grid/domain sensitivity & The current medium-to-fine comparison is preliminary only. \\
External reference tables & The benchmark here is the Ticket 01 European Black-Scholes call utility. \\
\bottomrule
\end{longtable}

\section{Lessons Learned / Study Notes}

The first lesson is that American option validation is not only about finding cases where early
exercise matters. It is also about checking cases where early exercise should not matter.

The second lesson is that \(U\geq\Phi\) is necessary but not sufficient. A no-dividend American
call can satisfy the obstacle and still be suspicious if it is materially above the European call.

The third lesson is that theorem checks require the theorem assumptions. Here, \(q=0\) is not a
minor parameter; it is the condition that makes the American and European call values equal.

The fourth lesson is that selected-spot rows are useful for human reading, while target-region
error metrics are better for numerical validation. Both are helpful, but they answer different
questions.

\section{Link to the Next Ticket}

The next call-validation ticket should move from the no-dividend control case to dividend-paying
American calls. With positive dividends, early exercise can become financially meaningful at high
spot prices. That future ticket should not reuse Ticket 06's equality conclusion; it should build
new validation criteria for the dividend-paying regime.

\section{Reproducibility Record}

\subsection{Files Created}

\begin{longtable}{P{0.54\textwidth}L{0.38\textwidth}}
\toprule
Path & Role in Ticket 06 \\
\midrule
\endhead
\file{experiments/03_no_dividend_american_call_validation.py} &
Runs the no-dividend American call validation cases and writes CSV/figure artifacts. \\
\file{tests/test_no_dividend_american_call_validation.py} &
Focused unittest coverage for Ticket 06 artifacts, metrics, scope, and docstrings. \\
\file{results/01_solver_validation/tables/ticket_06_no_dividend_american_call_validation.csv} &
Validation summary table. \\
\file{results/01_solver_validation/tables/ticket_06_no_dividend_american_call_selected_spots.csv} &
Selected moneyness price comparison table. \\
\file{results/01_solver_validation/figures/ticket_06_no_dividend_american_call_value_comparison.png} &
American call, European call, and payoff comparison figure. \\
\file{results/01_solver_validation/figures/ticket_06_no_dividend_american_call_theorem_error.png} &
American-minus-European theorem-error figure. \\
\file{reports/03_solver/tickets/ticket_06_no_dividend_american_call_validation.tex} &
Formal LaTeX study-report source. \\
\file{reports/03_solver/tickets/ticket_06_no_dividend_american_call_validation.pdf} &
Formal PDF compiled from the LaTeX source. \\
\bottomrule
\end{longtable}

\subsection{Commands}

Test command:
\begin{lstlisting}[language=bash]
PYTHONPYCACHEPREFIX=/private/tmp/codex_pycache PYTHONPATH=src MPLCONFIGDIR=/private/tmp/matplotlib-cache FC_CACHEDIR=/private/tmp/fontconfig-cache python3 -m unittest discover -s tests
\end{lstlisting}

Experiment command:
\begin{lstlisting}[language=bash]
PYTHONPYCACHEPREFIX=/private/tmp/codex_pycache PYTHONPATH=src MPLCONFIGDIR=/private/tmp/matplotlib-cache FC_CACHEDIR=/private/tmp/fontconfig-cache python3 experiments/03_no_dividend_american_call_validation.py
\end{lstlisting}

Compile command:
\begin{lstlisting}[language=bash]
PYTHONPYCACHEPREFIX=/private/tmp/codex_pycache PYTHONPATH=src python3 -m compileall -q src tests experiments
\end{lstlisting}

LaTeX compile command:
\begin{lstlisting}[language=bash]
HOME=/private/tmp FC_CACHEDIR=/private/tmp/fontconfig-cache latexmk -xelatex -interaction=nonstopmode -halt-on-error -outdir=/private/tmp/ticket06_latex reports/03_solver/tickets/ticket_06_no_dividend_american_call_validation.tex
\end{lstlisting}

PDF verification uses \code{pdfinfo} and page rendering with \code{pdftoppm}. Raw Git status is
intentionally not included in this formal PDF report.

\bibliographystyle{plainnat}
\bibliography{reports/03_solver/references}

\end{document}
